% ------------------------------------------------------------------------------
% Chapter 1
% ------------------------------------------------------------------------------
\chapter{Design Overview} 
\label{cha:design_overview}

\section{Controller}

A dual-master memory controller.
Master 0 (CPU) has fixed priority over Master 1 (DMA). Synchronous writes on clock, and combinational reads.

\subsection{Parameters}

There are two parameters to this dut:

\begin{itemize}
    \item \texttt{ADDR\_WIDTH} : address size.
    \item \texttt{DATA\_WIDTH} : data width.
\end{itemize}

\subsection{Functional Description}

\begin{table}[H]
    \centering
    \begin{tabular}{|l|l|}
        \hline
        \textbf{Condition}      & \textbf{Behavior}                                                                                 \\ \hline
        Only one master valid   & That master is serviced immediately                                                               \\ \hline
        Both masters valid      & Master 0 is serviced (priority). Master 1 is not granted until Master 0 de-asserts \texttt{req0}  \\ \hline
        No master valid         & Controller idle                                                                                   \\ \hline
    \end{tabular}
    \caption{Controller Behavior}
\end{table}

\begin{itemize}
    \item \textbf{Reads}: 0 clock cycles (combinational —immediate once address \& grant are stable).
    \item \textbf{Writes}: 1 clock cycle (synchronous —update happens on next rising edge when grant \& we are true).
\end{itemize}

\subsection{Interface}

\begin{table}[H]
    \centering
    \begin{tabular}{|l|c|c|l|}
        \hline
        \textbf{Signal}     & \textbf{Direction} & \textbf{Width} & \textbf{Description}            \\ \hline
        \texttt{clk}        & Input              & 1              & Clock signal                    \\ \hline
        \texttt{rst\_n}     & Input              & 1              & Active-low synchronous reset    \\ \hline
        \texttt{addr0}      & Input              & ADDR\_WIDTH    & Master 0 address                \\ \hline
        \texttt{wdata0}     & Input              & DATA\_WIDTH    & Master 0 data                   \\ \hline
        \texttt{rdata0}     & Output             & DATA\_WIDTH    & Master 0 data                   \\ \hline
        \texttt{we0}        & Input              & 1              & Master 0 write enable           \\ \hline
        \texttt{req0}       & Input              & 1              & Master 0 read request           \\ \hline
        \texttt{gnt0}       & Output             & 1              & Master 0 grant                  \\ \hline
        \texttt{addr1}      & Input              & ADDR\_WIDTH    & Master 1 address                \\ \hline
        \texttt{wdata1}     & Input              & DATA\_WIDTH    & Master 1 data                   \\ \hline
        \texttt{rdata1}     & Output             & DATA\_WIDTH    & Master 1 data                   \\ \hline
        \texttt{we1}        & Input              & 1              & Master 1 write enable           \\ \hline
        \texttt{req1}       & Input              & 1              & Master 1 read request           \\ \hline
        \texttt{gnt1}       & Output             & 1              & Master 1 grant                  \\ \hline
    \end{tabular}
    \caption{Controller signals}

\end{table}